% Appendix -- supporting material (<=4 pages for CoNEXT long paper)
\section*{Appendix}
\addcontentsline{toc}{section}{Appendix}

A.1 Coarse screen.  A candidate AS is admitted iff the
quick score computed from the week-over-week baseline comparison
clears the admission threshold.  The score aggregates a mean-level
shift and a dispersion ratio, each capped and weighted, so the
screen is monotone in the size of the deviation and independent of
the event's absolute traffic volume.  Candidates that fail are
marked \texttt{low-confidence} and routed to the degraded path.
\label{app:regularity}

A.2 Arbitration rules.  Six disambiguation rules operate
on structured detector arguments in fixed design-time
order: \texttt{leak\_induced\_hijack}, \texttt{multi-actor
leak}, \texttt{self-leak during outage},
\texttt{pure-leak + total-traffic-drop},
\texttt{ROA-vs-outage}, \texttt{leak-during-outage}.
The ordering is not tuned on the evaluation set.
\label{app:arbitration}

A.3 Abstention analysis.  The engine abstains at three
levels, and each level records its triggering statistic in
the cache, so abstentions are auditable rather than silent.
At the estimator level, the direction is committed only if
the bootstrap 95\% CI of that direction excludes zero and
the permutation $p$-value is below the significance
threshold; a leading series may also disqualify the
reverse direction outright.  At the synthesis level, the
engine abstains if the best channel's direction confidence
falls below the direction floor.  Finally, the Bayesian
synthesis reports \texttt{unknown} if the normalized
posterior of the dominant direction does not exceed the
commit floor $\tau = 0.34$.
\label{app:abstention}

A.4 Cache schema.  A versioned, content-addressed JSON store
with schema-version, deterministic hash, structured payload,
and input reference.  Every verdict, detector argument, and
abstention record is reproducible from this single artefact.
\label{app:cache}

A.6 Live-replay accounting.  The 192-alert live replay is
summarised into the evaluation cache
(\texttt{results/cloudflare\_alerts/run\_summary.json}),
which records the coverage rate, the per-source verdict
breakdown, and the per-alert timing distribution underlying
(L7).  Of 192 enqueued alerts, 130 completed adjudication and
62 failed in the harness; the 130 adjudicated alerts yield 14
\texttt{real\_incident} verdicts, split 7 curated and 7
discovered.\label{app:live-replay}


A.5 Confusion matrix.  Table~\ref{tab:dataset} reports
Top-1 and F1; Figure~\ref{fig:confusion} shows the
full attribution geometry.  The confusion matrix confirms
that mis-attributions cluster at class boundaries:
outages and route leaks share an AS-path signature,
and hijack victims are confused with route leaks when
multiple ASes announce the victim's prefix.
\label{app:confusion}

\begin{figure}[t]
\centering
\includegraphics[width=0.75\linewidth]{figs/fig_confusion.pdf}
\caption{Row-normalised attribution confusion matrix.
Diagonal cells show correct attribution; off-diagonal
cells show the specific mis-attribution geometry.
The 3 errors (1 outage, 1 leak, 1 hijack) all occur
at class boundaries.\label{fig:confusion}}
\end{figure}

A.7 Detector-output geometry for the 2025 events.
The 23 curated 2025 RRC00 events are projected into
the 5-D detector-output space (hijack, hijacking,
leak, outage, traffic) and reduced to two dimensions
by PCA. A month-by-cause heatmap shows events cluster
in time; an embedding map shows the four causes
separate cleanly with the three mis-attributed events
at cluster boundaries. Released as
\texttt{fig\_detector\_geometry.pdf}.
\label{app:case-geo}

A.8 Five-month deployment statistics.  Four panels:
rolling evaluation accuracy, daily raw-alarm trace
against BGP Tracer verdicts, monthly cause
distribution, and per-cause confidence histogram
showing bimodal profiles (hijack/leak concentrated,
outage dispersed). Released as
\texttt{fig\_deployment\_stats.pdf}.
\label{app:deployment}